\documentclass[11pt]{article}
\usepackage[margin=1.1in]{geometry}
\usepackage{amsmath,amssymb}
\usepackage{listings}
\usepackage{xcolor}
\usepackage[colorlinks=true,linkcolor=blue!50!black,urlcolor=blue!50!black]{hyperref}

\lstdefinestyle{lean}{
  basicstyle=\ttfamily\small,
  keywordstyle=\color{blue!60!black},
  commentstyle=\color{black!55},
  columns=fullflexible,
  breaklines=true,
  morekeywords={theorem,def,inductive,structure,Prop,Type,by,induction,cases}
}
\lstset{style=lean}
\newcommand{\lean}[1]{\texttt{\small #1}}

\title{Reagent: what the mechanised lemmas are for}
\author{}
\date{}

\begin{document}
\maketitle

\noindent
A plain-English guide to the Lean development: what each definition and lemma
exists to do, and why the proof needs it. Proofs are sketched in a line --- most
are inductions, and the content is \emph{which} induction and on \emph{what}.
Part~\ref{part:model} covers the model in \texttt{Reagent/Trace.lean};
Part~\ref{part:int} covers the proof of internal consistency in
\texttt{Reagent/IntConsistency.lean}.

\tableofcontents

\section{The model: \texttt{Trace.lean}}
\label{part:model}

\noindent
This note says, in plain English, what each lemma in \texttt{Trace.lean} exists
to do and why the development needs it. Proofs are only sketched --- almost
every one is an induction, and the interesting content is \emph{which} induction
and on \emph{what}.

\subsection{The setting, in one paragraph}

A trace is a list of labels. A label is only a description: a transaction that
reads the same location twice emits \lean{get(x,5)@t} twice, and those are two
different things that happen. So we work with \emph{events}: an event is
\lean{(id, label)}, where the \lean{id} is a \textbf{name}. The name has exactly
one job --- to tell two occurrences of one label apart. It is never compared,
never ordered, and it is not a position. What orders two events is the list they
sit in: \lean{OccursBefore es e f} holds when \lean{es} cuts as
\[
  \texttt{es} \;=\; p \;\mathbin{+\!\!+}\; [e] \;\mathbin{+\!\!+}\; m
  \;\mathbin{+\!\!+}\; [f] \;\mathbin{+\!\!+}\; s .
\]
Everything below is about turning that definition into a usable order, and about
getting from a trace (labels only) to a list of events (labels with names).

\subsection{The shape of the order}

\paragraph{\lean{OccursBefore.split} and \lean{OccursBefore.of\_split}.}
These are the two directions of the definition. \lean{split} takes a proof that
\lean{e} occurs before \lean{f} and hands back the cut $p, m, s$; \lean{of\_split}
takes a cut and hands back the proof. Between them, every later proof can stop
reasoning about the inductive rules \lean{base}/\lean{pre}/\lean{post} and reason
about lists instead, which is much easier. \lean{split} is by induction on the
derivation --- the three rules give the three cases --- and \lean{of\_split} is
the three constructors applied in order.

\subsection{What the names buy}

\paragraph{\lean{OccursBefore.id\_ne}.}
If \lean{e} occurs before \lean{f} in a list whose names are all distinct, then
\lean{e} and \lean{f} carry different names. Trivial-sounding, and it is the
whole reason events carry names at all.

\paragraph{\lean{OccursBefore.irrefl}.}
No event occurs before itself. This is \lean{id\_ne} read at $f := e$: it would
say $e.id \neq e.id$. Note what fails without names: on a trace
\lean{begin@t} $\cdot$ \lean{get(x,5)@t} $\cdot$ \lean{get(x,5)@t}, a relation on
\emph{labels} really does put \lean{get(x,5)@t} before \lean{get(x,5)@t}, so it
is reflexive and cannot be a strict order. The names are what make those two
occurrences two events, and irreflexivity is the first thing they pay for.

\paragraph{\lean{UniqueIds.split\_unique}.}
If a list cuts around one event in two different ways, the two cuts are the same
cut. This is the precise sense in which a name \emph{locates} an event: it occurs
in one place. Proof is an induction on the first prefix, using distinctness to
rule out the mismatched cases.

\paragraph{\lean{OccursBefore.trans}.}
If $e$ is before $f$ and $f$ is before $g$, then $e$ is before $g$. The proof is
the reason \lean{split\_unique} exists: take the cut around $e,f$ and the cut
around $f,g$; both are cuts of the same list at the same event $f$, so they line
up, and gluing them puts $f$ inside the middle segment of a single cut around
$e,g$. Nothing here compares names as numbers --- distinctness is used only to
know that the two cuts are cuts \emph{at the same place}.

\subsection{Totality}

\paragraph{\lean{mem\_split}.}
An element of a list splits the list around itself. Pure list plumbing; Lean's
core does not ship it in this form.

\paragraph{\lean{OccursBefore.total}.}
Any two events of one list are ordered one way or the other, or are the same
event. Worth noticing: this needs \emph{no} hypothesis about names. Membership
alone decides it, because a list already puts one of its elements before the
other. Names are needed for irreflexivity and transitivity, not for totality.

\subsection{Program order}

\lean{po} of a transaction is defined, not assumed: \lean{poOf es t} is just
\lean{OccursBefore es} restricted to that transaction's events (the restriction
is carried by the subtype, since \lean{eventsOf es t} already demands the tag).
So its order properties are theorems, and Cerone's model requires $(E, po)$ with
$po$ a total order --- that is the obligation these four discharge.

\begin{itemize}
  \item \lean{poOf\_irrefl}, \lean{poOf\_trans} --- inherited directly from
        \lean{OccursBefore}.
  \item \lean{poOf\_antisymm} --- holds because \lean{po} is \emph{strict}: its
        two hypotheses together give $a$ before $a$, which is impossible. If one
        prefers the reflexive reading of $po$, this is the clause that stops
        being vacuous and starts doing work.
  \item \lean{poOf\_total} --- trichotomy on the transaction's events.
\end{itemize}

All but \lean{poOf\_total} take \lean{UniqueIds es} as a hypothesis, which is
discharged once and for all in Section~\ref{sec:naming}.

\subsection{From a trace to events: one step at a time}

The operational semantics produces a run, and a run grows at the \emph{back} ---
each step appends its label to what has been emitted. Lists, however, recurse at
the front. So proofs about runs need snoc-shaped facts: what does \lean{events}
do when one more label arrives?

\paragraph{\lean{helperEventsFrom\_append}.}
Naming $\tau \mathbin{+\!\!+} \sigma$ is naming $\tau$, then naming $\sigma$ with
the counter started at $i + |\tau|$. The general statement; the next two are the
cases actually used.

\paragraph{\lean{events\_snoc\_ev} and \lean{events\_snoc\_adm}.}
One more transition label appends exactly one event at the end;
one more administrative step appends nothing. These two are what let an induction
over a run proceed: at each step you know precisely how the event list changed,
so an invariant stated about \lean{events $\tau$} can be pushed through the step.
Every case of \lean{inv\_store} and \lean{inv\_int} in
\texttt{IntConsistency.lean} begins by rewriting with one of them.

\paragraph{\lean{helper\_id\_ge}, \lean{helper\_id\_lt}, \lean{events\_id\_lt}.}
Every name handed out while naming $\tau$ lies below $|\tau|$. The point is
\emph{freshness}: the event appended by the next step is named $|\tau|$, so it is
not a name already in use. This is the only place the counter's increasing-ness
is used, and note what it is used \emph{for} --- to show names do not repeat,
never to conclude that one event precedes another.

\subsection{The naming, discharged}
\label{sec:naming}

\paragraph{\lean{uniqueIds\_helper} and \lean{uniqueIds\_events}.}
\lean{UniqueIds (events $\tau$)}: the canonical naming really does hand out
distinct names. This is the hypothesis every \lean{poOf\_*} theorem carries, so
proving it once means the consistency proofs never have to drag it around.

\paragraph{\lean{map\_helper} and \lean{eventsOf\_events}.}
\lean{EventsOf $\tau$ (events $\tau$)}: the canonical naming satisfies the
\emph{specification} of a naming --- its labels spell out $\tau$ exactly, in
order, with the administrative steps dropped, and its names are distinct.

This pair deserves a comment, because it is the only place where the two halves
of the design meet. \lean{EventsOf} is a predicate: it says what it \emph{takes}
to be a naming of $\tau$, without picking one, so any list of events meeting it
would do --- the names may be $7, 3, 99$ for all it cares. \lean{events} is a
function: it picks one, because extracting transactions from a trace has to
produce something concrete. \lean{eventsOf\_events} is the proof that the choice
is legitimate. Without it, \lean{events} would be an assumption smuggled in as a
definition; with it, the choice is a witness, and every theorem stated for an
arbitrary naming applies to it.

\subsection{Where they are used}

\begin{itemize}
  \item \lean{split}, \lean{of\_split}, \lean{split\_unique} --- machinery; used
        only to prove the next two.
  \item \lean{irrefl}, \lean{trans} --- used constantly, and in particular to
        show that an event cannot be both before and after another, which is how
        maximality arguments (``the last access before this read'') close.
  \item \lean{total} --- what makes $(E, po)$ a total order, as the declarative
        model demands.
  \item \lean{events\_snoc\_ev}, \lean{events\_snoc\_adm} --- the spine of every
        induction over a run: \lean{inv\_nextTx}, \lean{inv\_store},
        \lean{inv\_int}.
  \item \lean{events\_id\_lt}, \lean{uniqueIds\_events} --- freshness and
        distinctness, i.e.\ the hypotheses everything else assumes.
  \item \lean{eventsOf\_events} --- non-vacuity of the whole construction.
\end{itemize}

\section{Internal consistency: \texttt{IntConsistency.lean}}
\label{part:int}

\noindent
Part~\ref{part:model} covered the model; this part covers the proof of internal
consistency built on it. What matters here is which object each lemma bridges to
which.

\subsection{The gap the file has to close}

\lean{Int} is a statement about \emph{events}: a read must carry the value of the
$po$-last access to $x$ in its own transaction. The semantics knows nothing about
events. Its rule for reading is
\[
  \sigma(l) = v \;\Longrightarrow\;
  \langle \sigma, i, t, \ldots\rangle \;\xrightarrow{\ \mathrm{get}(l,v)@t\ }\;
  \langle \sigma, i, t, \ldots\rangle,
\]
that is: the value returned is whatever the \emph{store} holds, and the rule says
nothing about what the transaction did earlier. So the whole file is one bridge:
\begin{center}
  \emph{why is the store's current value the one this transaction last left there?}
\end{center}
Under the lock it is, and that fact is \lean{inv\_store}. Everything before it is
machinery to state and prove it; everything after it carries the fact up to the
declarative model.

\subsection{The axiom, restated}

\paragraph{\lean{Label.accesses} and \lean{Label.accessVal}.}
$HEvent_x$ --- is this label a read or a write on $x$ --- and the value such a
label carries. \lean{accessVal} is an \lean{Option}, \lean{none} on
\lean{begin}/\lean{eend}/\lean{retry}, which is how the type records that those
carry no value.

\paragraph{\lean{Transaction.PriorAccess} and \lean{MaxPriorAccess}.}
$po^{-1}(e) \cap HEvent_x$, and its $po$-greatest element. Because $po$ is strict
here, \lean{PriorAccess} does not have to say $a \neq e$ --- the old
formulation did, when $po$ was $\leq$.

\paragraph{\lean{Transaction.Int} and \lean{History.Int}.}
The axiom verbatim for one transaction, then quantified over the history's
transactions. Note what is absent: \lean{ar} never appears. \lean{Int} is
intra-transactional, which is why it can be proved for \emph{every} history of a
run rather than for some chosen one.

\subsection{Inducting over a run}

\paragraph{\lean{Run.split}.}
A run splits wherever its trace does: if the trace is $\mathit{pre} \mathbin{+\!\!+}
\mathit{post}$, there is a middle configuration.

\paragraph{\lean{listReverseInduction} and \lean{Run.reverseRecOn}.}
A run grows at the back, so the useful induction is ``a run that emitted $\tau$,
then took one more step''. Lean's list recursor goes the other way, so
\lean{listReverseInduction} builds the snoc recursor by running the ordinary one
over the reversed list, and \lean{Run.reverseRecOn} lifts that to runs. Every
invariant below is proved with it; without it each proof would have to reason
about a step at the \emph{front}, where the configuration is not yet the one the
invariant talks about.

\paragraph{\lean{OccursBefore.snoc\_last} and \lean{snoc\_mono}.}
The two facts about appending one event: everything already in the list occurs
before the newcomer, and an order that already held still holds. The second is
just the \lean{post} constructor. These are what let a maximality claim survive a
step --- the newcomer is after everything, so it cannot sneak between two older
events.

\subsection{The invariants}

\paragraph{\lean{inv\_nextTx}.}
Every tag occurring in the run, and the live tag, is below \lean{nextTx}. This is
the mechanised form of ``\lean{begin} mints a \emph{fresh} id''. Its purpose is
one case of the next proof: when a transaction begins, the invariant about its
past has to hold vacuously, and that is exactly the claim that nothing in the run
so far carries its tag. Without freshness a reused id would splice two unrelated
attempts into one transaction.

\paragraph{\lean{LastAccess}.}
``$a$ is the last thing $t$ did to $x$ so far'': $a$ is an access of $t$ to $x$,
and no access of $t$ to $x$ occurs after it. Stated with \lean{OccursBefore},
not by comparing names.

\paragraph{\lean{inv\_store} --- the heart of the file.}
If $t$ holds the lock, then for every location $x$, the value in the store is the
one carried by $t$'s last access to $x$. Proved by induction over the run, one
case per rule, and the lock discipline is spent in exactly three of them:
\begin{itemize}
  \item \lean{get}/\lean{set} on $x$: the new event \emph{is} the last access ---
        anything else would be before it, contradicting maximality --- and the
        value it carries is what the rule put in the store.
  \item \lean{begin}: the tag is fresh, so by \lean{inv\_nextTx} the hypothesis
        ``this tag already has an access'' is false and the case is vacuous. A
        transaction starts with no history.
  \item \lean{eend}, \lean{retry}: both clear \lean{activeTx} in the same step,
        so there is no live transaction for the invariant to constrain. This is
        why \lean{retry}'s rollback of the store is harmless: a rollback never
        lands inside a live transaction.
\end{itemize}

\subsection{From the store back to the events}

\paragraph{\lean{OccursBefore.mem}.}
Related events are elements of the list. Plumbing.

\paragraph{\lean{OccursBefore.restrict\_append}.}
If $a$ and $b$ both lie in $A$, and $a$ occurs before $b$ in $A \mathbin{+\!\!+} B$,
then $a$ occurs before $b$ in $A$ alone. This is the lemma that lets an older
read's answer be computed in the prefix of the run rather than in the whole run
--- without it, ``nothing lies between $a$ and $e$'' could not be transported
from the long list to the short one. It needs distinct names: with duplicates,
$b$ could occur twice and the cut found in $A \mathbin{+\!\!+} B$ need not be a cut
of $A$. The proof splits $A$ at $a$ and asks where $b$ sits; after $a$ it is a
direct cut, before $a$ it contradicts transitivity.

\paragraph{\lean{uniqueIds\_helper}, \lean{uniqueIds\_events}.}
\lean{UniqueIds (events $\tau$)} --- the hypothesis that \lean{irrefl},
\lean{trans} and \lean{restrict\_append} all carry, discharged for the canonical
naming so the main proofs never mention it.

\paragraph{\lean{isAccess\_of\_accesses}.}
Accessing $x$ makes a label a \lean{get} or a \lean{set}. Small, and needed at one
precise point: to place an event into a transaction's $E$, whose membership test
demands \lean{isAccess}.

\paragraph{\lean{LastAccessBefore}, \lean{not\_occursBefore\_snoc},
\lean{lastAccessBefore\_of\_snoc}.}
``$a$ is the last access of $t$ to $x$ \emph{before} $e$'' --- the relativised
form of \lean{LastAccess}, because \lean{Int} speaks about the state of affairs
at a particular read, not at the end of the run. The other two say that appending
a step does not disturb an older read's answer: nothing occurs after the newcomer,
so the newcomer cannot be the thing that was missing.

\subsection{The two theorems}

\paragraph{\lean{inv\_int}.}
\lean{Int} on the canonical naming: every read sees the value its own transaction
last left at that location. Induction over the run again. Every case is
bookkeeping except one --- a \lean{get} that \emph{is} the newly appended event,
where \lean{inv\_store} on the prefix says the store already holds what the
transaction last left, and the rule says the read returns the store. That single
case is where the whole file lands.

\paragraph{\lean{int\_of\_run}.}
The translation from the canonical naming to an arbitrary history of the trace.
\lean{HistoryOf} forces $T = \lean{transactionOf (events $\tau$) $t_0$}$, so $T.E$
and $T.po$ unfold to the definitions, and \lean{MaxPriorAccess} becomes
\lean{LastAccessBefore}. The only real content is maximality, and it is
contradiction-shaped: given some $b$ lying between $a$ and the read, wrap $b$ as
an element of $T.E$, feed it to the maximality hypothesis, and get either $b = a$
or $b$ before $a$ --- both contradict $a$ before $b$ by irreflexivity and
transitivity. Notably \lean{poOf\_total} is \emph{not} needed: nothing has to
compare two arbitrary events, only to rule out an impossible one.

\subsection{What is proved, and what is still assumed}

With \lean{int\_of\_run} closed, internal consistency holds for every run of the
semantics and \emph{every} history of that run:
\[
  \forall \tau, \forall H. \quad
  \mathrm{Run}\;C_{\mathrm{init}}\;\tau\;C' \;\wedge\; \mathrm{HistoryOf}\;\tau\;H
  \;\Longrightarrow\; H \models \mathrm{Int}.
\]
Two caveats worth keeping in view. First, the theorem says nothing if no $H$
satisfies \lean{HistoryOf $\tau$ H}; that is what \lean{exists\_historyOf} is for,
and it is still open. Second, \lean{Ext} is a different matter entirely: it reads
\lean{ar}, which is assumed rather than defined, so it cannot be proved for every
history once transactions may overlap --- under the lock \lean{ar} is forced, and
the $\forall H$ form survives.

\end{document}
